\subsection{GAP による群作用の確認}

\subsubsection{目的}

有限群の群作用について、GAP を用いて具体的な計算を行う。

特に、群作用における軌道と安定化群を計算し、
軌道安定化群定理
\[
|Gx|=[G:G_x]
\]
を確認する。

ここで、$G$ は有限群、$X$ は $G$ が作用する集合、
$x\in X$ に対して
\[
Gx=\{gx\mid g\in G\}
\]
を $x$ の軌道、
\[
G_x=\{g\in G\mid gx=x\}
\]
を $x$ の安定化群とする。

\subsubsection{数学的背景}

群 $G$ が集合 $X$ に作用しているとする。

$x\in X$ に対して、その軌道は
\[
Gx=\{gx\mid g\in G\}
\]
であり、安定化群は
\[
G_x=\{g\in G\mid gx=x\}
\]
である。

軌道安定化群定理によれば、
\[
|Gx|=[G:G_x]
\]
が成り立つ。

さらにラグランジュの定理より、
\[
[G:G_x]|G_x|=|G|
\]
なので、
\[
|Gx||G_x|=|G|
\]
を得る。

これは、前の例で確認したラグランジュの定理と直接つながっている。

\subsubsection{$S_3$ の自然な作用}

具体例として、3文字の対称群
\[
S_3
\]
を考える。

$S_3$ は集合
\[
X=\{1,2,3\}
\]
に自然に作用する。

例えば、置換
\[
(1\,2)
\]
は
\[
1\mapsto2,\qquad
2\mapsto1,\qquad
3\mapsto3
\]
と作用する。

この作用において $1$ の軌道を考えると、
\[
S_3\cdot1=\{1,2,3\}
\]
となる。

一方、$1$ の安定化群は
\[
(S_3)_1
=
\{g\in S_3\mid g(1)=1\}
\]
である。

これは
\[
(S_3)_1=\{(),(2\,3)\}
\]
なので、
\[
|(S_3)_1|=2
\]
である。

したがって、
\[
|S_3\cdot1|=3,\qquad
|(S_3)_1|=2
\]
であり、
\[
|S_3\cdot1|\cdot|(S_3)_1|
=3\cdot2
=6
=|S_3|
\]
となる。

\subsubsection{GAP による計算}

以下のプログラムを用いる。

\begin{lstlisting}[language={}]
G := SymmetricGroup(3);
X := [1..3];

orbit := Orbit(G, 1);
stab := Stabilizer(G, 1);

Size(G);
orbit;
Size(orbit);
stab;
Size(stab);
Index(G, stab);
Size(orbit) * Size(stab);
\end{lstlisting}

実行すると、概ね次のような結果が得られる。

\begin{lstlisting}
Group([ (1,2), (1,2,3) ])
[ 1, 2, 3 ]
3
Group([ (2,3) ])
2
3
6
\end{lstlisting}

ここから、

\[
|G|=6,
\]
\[
G\cdot1=\{1,2,3\},
\]
\[
|G\cdot1|=3,
\]
\[
|G_1|=2,
\]
\[
[G:G_1]=3
\]

であることが分かる。

したがって、

\[
|G\cdot1|
=
[G:G_1]
=
3
\]

となり、軌道安定化群定理が確認できる。

さらに、

\[
|G\cdot1||G_1|
=
3\cdot2
=
6
=
|G|
\]

となる。

\subsubsection{すべての軌道を調べる}

GAP では、ある集合上のすべての軌道を
\verb|Orbits| を使って調べることもできる。

\begin{lstlisting}[language={}]
Orbits(G, X);
\end{lstlisting}

この場合、$S_3$ は $\{1,2,3\}$ 上推移的に作用しているので、

\begin{lstlisting}
[ [ 1, 2, 3 ] ]
\end{lstlisting}

が得られる。

つまり、$X$ 全体が一つの軌道になっている。

これは
\[
S_3\cdot1
=
S_3\cdot2
=
S_3\cdot3
=
X
\]
であることを意味する。

% \subsubsection{ラグランジュの定理との関係}

% 前の例では、部分群 $H\leq G$ に対して
% \[
% |G|=[G:H]|H|
% \]
% を確認した。

% 群作用の場合、点 $x$ の安定化群 $G_x$ は $G$ の部分群なので、
% ラグランジュの定理を適用することができる。

% 特に、
% \[
% |G|=[G:G_x]|G_x|
% \]
% である。

% 軌道安定化群定理
% \[
% |Gx|=[G:G_x]
% \]
% を使うと、
% \[
% |G|=|Gx||G_x|
% \]
% が得られる。

% したがって、群作用を考えることで、
% ラグランジュの定理をより一般的な構造の中で見ることができる。

\subsubsection{確認したこと}

今回の計算では、GAP を用いて

\begin{itemize}
  \item 群の軌道を計算する。
  \item 点の安定化群を計算する。
  \item 軌道の大きさを計算する。
  \item 安定化群の指数を計算する。
\end{itemize}

ことができた。

具体的に $S_3$ の自然な作用について、

\[
|S_3\cdot1|
=
[S_3:(S_3)_1]
=
3
\]

を確認した。

さらに、

\[
|S_3\cdot1|\cdot|(S_3)_1|
=
|S_3|
\]

も計算によって確認できた。

これは、ラグランジュの定理から軌道安定化群定理へ進む具体例になっている。

\subsubsection{使用したもの}

\begin{itemize}
  \item GAP
  \item 有限群論
  \item 群作用
  \item 軌道
  \item 安定化群
  \item 軌道安定化群定理
\end{itemize}